\bookchapter{פרק 33: מבט לעתיד – הצעד הבא של הטלפון}

\chapterintro{לא יעלם, אך יאבד מונופול. 3 נושאים: מכשירים, “היעלמות”, אתיקה.}

\booksection{א', AR/VR: “שולחן שני”}

התחרות החלה. AR/VR:

\begin{itemize}
  \item \textbf{Vision Pro}:Apple Vision Pro = "מרחב". מחיר/משקל. 2026-2027 זול/קל. iPhone = מרכז נתונים.
  \item \textbf{ממשק מוח-מחשב}:"שולחן שני". Meta Ray-Ban, Rokid, XREAL, INMO. 2025-2026: תרזום/HUD/AI. 3-5 שנים להחלפת טלפון.
  \item \textbf{טבעות/לבישים}:Galaxy Ring, Oura. טלפון = מרכז + חלקים.
\end{itemize}

\booksection{ב', BCI}

מחשבה כממשק, אך ארוך:

\begin{itemize}
  \item \textbf{Neuralink}:2024 השתלה ראשונה, עכבר. למשותקים.
  \item \textbf{BCI לא-פולשני}:Kernel, OpenBCI, ללא ניתוח, אך פחות מדויק.
  \item \textbf{צרכני}:10-20 שנה. נגישות/ריכוז. אתיקה/פרטיות.
\end{itemize}

\booksection{ג', “היעלמות”}

כיוון: “היעלמות” - זהות דיגיטלית:

\begin{itemize}
  \item \textbf{מורשת דיגיטלית כהגדרת יצרן}:אימות מיידי.
  \item \textbf{מ"מרכז מכשירים" ל"מרכז זהות"}:ירושה כברירת מחדל.
  \item \textbf{שלוש. הרהור אתי: האם אנחנו עדיין הבעלים של הטלפון?}:ענן = מרכז, חומרה = כניסה.
\end{itemize}

\booksection{ד', אתיקה}

כשאתה מדליק את המסך בלי משים כל כמה דקות, זה הרגל שלך או “התנהגות ממכרת” מתוכננת?

\begin{itemize}
  \item \textbf{אתיקת AI}:AI > אתה? אחריות?
  \item \textbf{ריבונות נתונים}:מיינך? GDPR = התחלה.
  \item \textbf{פער דיגיטלי}:AI/מתקפל/לוויין = פער. כמו "לגלוש" לפני 20 שנה.
  \item ארבע. להישאר מחוברים, להישאר עצמאיים – פילוסופיית החיים הדיגיטליים שלך
  \item היו מקצי משאבים, לא מקבלים פסיביים של תוכן נדחף.
\end{itemize}

\booksection{ה', פילוסופיה}

10 שנה:

\begin{enumerate}
  \item בנו "עורף דיגיטלי" יציב.
  \item בצעו באופן קבוע "ניקוי דיגיטלי" ו"צום חברתי".
  \item תמיד הציבו אבטחה ופרטיות כתנאי מקדים לכל נוחות.
  \item בואו ננסה: נסחו את "תנאי השימוש האישיים שלכם למכשיר דיגיטלי"
  \item כתבו לכם ולטלפון שלכם "תנאי שימוש" פשוטים, שיגדירו את מערכת היחסים ביניכם. לדוגמה:
\end{enumerate}

\begin{tcbpractice}{סעיף 1: זמן הקשב שלי שייך לי. כשאני עובד/מבלה עם המשפחה, עליך להיכנס באופן יזום למצב "נא לא להפריע".}
“תנאים” לטלפון שלך:

\begin{itemize}
  \item סעיף 3: המטרה המשותפת שלנו היא לשפר את יעילות ואיכות החיים. אם תגרום לי להתמכר באמצעות עיצוב מניפולטיבי, שמורה לי הזכות להסיר אותך.
  \item חמש. נספחים
  \item ובכן, אחרי המסע הארוך של יותר מעשרים פרקים, כל הכבוד! התקדמתם בהצלחה מ"טירון טלפונים" ושלטתם במיומנויות הליבה ממתחילים ועד מתקדמים. נספח זה הוא כמו מדריך כיס ומאגר משאבים עבורכם.
\end{itemize}

\end{tcbpractice}

\begin{tcbsummary}{סיכום פרק}
\begin{itemize}
  \item Vision/AR = כניסה חדשה, לא מחליף
  \item BCI 10-20 שנה, אתיקה קודם
  \item טלפון "היעלמות" זהות
  \item AI/ריבונות/פער = 10 שנה
  \item חיבור + עצמאות
\end{itemize}

\end{tcbsummary}
